%% LyX 2.4.4 created this file.  For more info, see https://www.lyx.org/.
%% Do not edit unless you really know what you are doing.
\documentclass[english]{beamer}
\usepackage{mathptmx}
\usepackage[T1]{fontenc}
\usepackage[latin9]{inputenc}
\usepackage{units}
\usepackage{amssymb}
\usepackage{graphicx}

\makeatletter

%%%%%%%%%%%%%%%%%%%%%%%%%%%%%% LyX specific LaTeX commands.
%% A simple dot to overcome graphicx limitations
\newcommand{\lyxdot}{.}


%%%%%%%%%%%%%%%%%%%%%%%%%%%%%% Textclass specific LaTeX commands.
% this default might be overridden by plain title style
\newcommand\makebeamertitle{\frame{\maketitle}}%
% (ERT) argument for the TOC
\AtBeginDocument{%
  \let\origtableofcontents=\tableofcontents
  \def\tableofcontents{\@ifnextchar[{\origtableofcontents}{\gobbletableofcontents}}
  \def\gobbletableofcontents#1{\origtableofcontents}
}

%%%%%%%%%%%%%%%%%%%%%%%%%%%%%% User specified LaTeX commands.
\usetheme{Warsaw}
% or ...

\setbeamercovered{transparent}
% or whatever (possibly just delete it)
\setbeamercovered{invisible} %makes overlays invisible


%gets rid of navigation symbols
\setbeamertemplate{navigation symbols}{}


\usepackage{pgfpages}
%\pgfpagesuselayout{4 on 1}[a4paper, landscape, border shrink=7mm]
%\pgfpageslogicalpageoptions{1}{border code=\pgfusepath{stroke}}
%\pgfpageslogicalpageoptions{2}{border code=\pgfusepath{stroke}}
%\pgfpageslogicalpageoptions{3}{border code=\pgfusepath{stroke}}
%\pgfpageslogicalpageoptions{4}{border code=\pgfusepath{stroke}}

\makeatother

\usepackage{babel}
\begin{document}
\title[Short Title]{9.4: Divergence and Integral Tests}
\subtitle{Calculus III | MTH 331~~|~~Section A~~|~~Fall 2012}
\author{Dr.\,James Ashe}
\titlegraphic{\includegraphics[height=1.5cm]{D:/home/jim/JamesAshe/JCSU/jcsu_bull.jpg}}
\institute{Johnson C. Smith University}
\date{{}}

\makebeamertitle

%\beamerdefaultoverlayspecification{<+->}
\begin{frame}{}

\begin{block}{The Harmonic Series}

${\displaystyle \sum_{k=1}^{\infty}\frac{1}{k}=1+\frac{1}{2}+\frac{1}{3}+\frac{1}{4}+\dots}$
\end{block}
\begin{itemize}
\item <2->[]$S_{1}=1$
\item <3->[]$S_{2}=1+\frac{1}{2}=\frac{3}{2}$
\item <4->[]$S_{3}=1+\frac{1}{2}+\frac{1}{3}=\frac{11}{6}$
\item <5->[]$\vdots$
\item <6->[]$S_{10^{40}}\approx92.68$ \uncover<7->{~;$10^{40}\geq$drops
of water in the Pacific Ocean.}
\end{itemize}
\end{frame}

\begin{frame}{Properties of Convergent Series}

\begin{itemize}
\item [ex.]${\displaystyle \sum_{k=1}^{\infty}\mbox{\ensuremath{\left[3\left(\frac{1}{2}\right)^{k}+5\left(\frac{3}{4}\right)^{k}\right]}}}$\uncover<2->{$={\displaystyle \sum_{k=1}^{\infty}3\left(\frac{1}{2}\right)^{k}+\sum_{k=1}^{\infty}5\left(\frac{3}{4}\right)^{2}}$}
\textcolor{blue}{}\\
\textcolor{blue}{\uncover<3->{$\sum(a_{k}\pm b_{k})=\sum a_{k}\pm\sum b_{k}$}}
\item <4->[]$={\displaystyle \sum_{k=1}^{\infty}3\left(\frac{1}{2}\right)^{k}+\sum_{k=1}^{\infty}5\left(\frac{3}{4}\right)^{k}}$
\uncover<5->{$={\displaystyle 3\sum_{k=1}^{\infty}\left(\frac{1}{2}\right)^{k}+5\sum_{k=1}^{\infty}\left(\frac{3}{4}\right)^{k}}$}
\textcolor{blue}{}\\
\textcolor{blue}{\uncover<6->{$\sum c\cdot a_{k}=c\sum a_{k}$}}
\item <7->[]$=3\cdot\frac{\nicefrac{1}{2}}{1-\nicefrac{1}{2}}+5\cdot\frac{\nicefrac{3}{4}}{1-\nicefrac{3}{4}}=3\cdot1+5\cdot3=18$
\item <8->Also, if ${\displaystyle \sum_{k=0}^{\infty}a_{k}}$ converges/diverges
than ${\displaystyle \sum_{k=m}^{\infty}a_{k}}$ also converges/diverges.
\end{itemize}
\end{frame}

\begin{frame}{The Divergence Test}

\begin{block}{Theorem: Divergence Test}

\textbf{a. }If ${\displaystyle \sum a_{k}}$ converges, then ${\displaystyle \{a_{k}\}^{\infty}=0}$
(${\displaystyle \lim_{k\rightarrow\infty}a_{k}=0}$). \textbf{}\\
\textbf{b. }If $\{a_{k}\}^{\infty}\neq0$ then $\sum a_{k}$ diverges.
\end{block}
\begin{itemize}
\item <2->[ex.]${\displaystyle \lim_{k\rightarrow\infty}\frac{2k}{2k+1}}={\displaystyle \lim_{k\rightarrow\infty}\frac{\nicefrac{2k}{k}}{\nicefrac{2k}{k}+\nicefrac{1}{k}}={\displaystyle \lim_{k\rightarrow\infty}\frac{2}{2}=1}}$\uncover<3->{$\Rightarrow$${\displaystyle \sum_{k=2}^{\infty}\frac{2k}{2k+1}}$}\uncover<4->{
diverges.}
\item <5->[ex.]${\displaystyle \lim_{k\rightarrow\infty}\frac{1}{\sqrt{k}}=0}$
\uncover<6->{, but we \textbf{cannot} say that ${\displaystyle \sum_{k=1}^{\infty}\frac{1}{\sqrt{k}}}$
converges; the Divergence test is inconclusive.}
\item <7->The Divergence Test allows us to evaluate a series by looking
at a sequence.
\end{itemize}
\end{frame}

\begin{frame}{The Integral Test}

\begin{block}{Theorem: Integral Test}

If $f$ is a continuous, positive, decreasing function for $x\geq1$
such that $a_{k}=f(k)$ for $k=1,2,\dots$ then ${\displaystyle \sum_{k=1}^{\infty}a_{k}}$
and $\int_{1}^{\infty}f(x)dx$ either \textbf{both }converge or diverge.
\end{block}
\begin{itemize}
\item <2->If they converge they do \textbf{not }have to converge to the
same value.
\item <3->[]\includegraphics[width=5cm]{images/300px-Integral_Test\lyxdot svg}
\end{itemize}
\end{frame}

\begin{frame}{}

\begin{itemize}
\item [ex.]Show that the Harmonic Series, ${\displaystyle \sum_{k=1}^{\infty}\frac{1}{k}}$
diverges.
\item []<2->Is $f(x)=\frac{1}{x}$ continuous for $x\geq1$? \textcolor{green}{\uncover<3->{\checkmark}}
\item []<4->Is $f(x)$ positive for $x\geq1$? \textcolor{green}{\uncover<5->{\checkmark}}
\item []<6->Is $f(x)$ decreasing, i.e., \textcolor{black}{$f'(x)=\frac{d}{dx}f(x)<0$,
for $x\geq1$?} \textcolor{green}{\uncover<7->{}}
\item []<7-> \hspace{.3cm}$f'(x)=-x^{-2}<0$\textcolor{green}{{} \checkmark}
\item []<8->So we can apply the Integral Test:
\item []<9->\hspace{.3cm}$\int_{1}^{\infty}\frac{1}{x}dx={\displaystyle \lim_{b\rightarrow\infty}\ln x\vert_{1}^{b}=}\lim_{b\rightarrow\infty}\left[\ln(b)-\ln(1)\right]=$\uncover<11->{$\infty$}
\end{itemize}
\end{frame}

\begin{frame}{}

Recall: Integration by substitution, 
\begin{itemize}
\item [ex.]Find$\int\frac{1}{x\ln x}dx$. \uncover<2->{Let $u=\ln x$}
\item <3->[]$\frac{du}{dx}=\frac{1}{x}\Rightarrow dx=x\,du$
\item <4->[]So $\int\frac{1}{x\ln x}dx=\int\frac{1}{x\cdot u}\cdot x\,du=\int\frac{1}{u}\,du=\ln u=\ln(\ln x)$
\item []
\item [ex.\,23]Does ${\displaystyle \sum_{k=2}^{\infty}\frac{1}{k\ln k}}$
converge or diverge?
\item <5->[]First note that $f(x)=\frac{1}{x\ln x}$ is decreasing (check:
$\frac{1}{x\ln x}\frac{d}{dx}<0)$, positive, and continuous for $x\geq1$.
\item <6->[]$\int_{2}^{\infty}\frac{1}{x\ln x}dx={\displaystyle \lim_{b\rightarrow\infty}}\ln(\ln x)\vert_{2}^{b}={\displaystyle \lim_{b\rightarrow\infty}\left[\ln(\ln b)-\ln(\ln2)\right]=\infty}$
\end{itemize}
\end{frame}

\begin{frame}{Estimating Infinite Series with Integrals}

While in general when $\int a_{k}$ and $\sum a_{k}$ converge it
is to different values, we can \emph{estimate }$\sum a_{k}$ using
$\int a_{k}$.
\begin{block}<2->{Theorem: Estimating Series with Positive Terms}

Let $f$ is positive, decreasing, and continuous for $x\geq1$ and
$f(k)=a_{k}$ for $k=1,2,\dots$. If ${\displaystyle \sum_{k=1}^{\infty}a_{k}}$converges
to $S$ then, $S_{n}+\int_{n+1}^{\infty}f(x)dx{\displaystyle \leq\sum_{k=1}^{\infty}a_{k}}\leq S_{n}+\int_{n}^{\infty}f(x)dx$. 
\end{block}
\begin{itemize}
\item <3->$R_{n}=S-S_{n}$ is called the \textbf{remainder}.
\end{itemize}
\begin{columns}


\column{5cm}
\begin{itemize}
\item []<3-|handout:0>\includegraphics[width=3cm]{images/int_estimate1}
\item []<3->$\int_{n+1}^{\infty}f(x)dx\leq R_{n}$
\end{itemize}

\column{5cm}
\begin{itemize}
\item []<4-|handout:0>\includegraphics[width=3cm]{images/int_estimate2}
\item []<4->$R_{n}\leq\int_{n}^{\infty}f(x)dx$
\end{itemize}
\end{columns}
\end{frame}

\begin{frame}{}

\begin{itemize}
\item [ex.]Using $n=10$, estimate the value of ${\displaystyle \sum_{n=1}^{\infty}\frac{1}{n^{2}}}$
\\
(it is known to be $\frac{\pi^{2}}{6}\approx1.64493$.)
\item []<2->First get the partial sum, $S_{10}\approx1.54977$.
\item []<3->Next compute the integrals, for $n=10$: $\int_{10}^{\infty}\frac{1}{x^{2}}dx=\frac{1}{10}$
and for $n+1=11$: $\int_{11}^{\infty}\frac{1}{x^{2}}dx=\frac{1}{11}$
\item []<4->Using the Theorem:\uncover<5->{ $S_{10}+\frac{1}{11}\leq S\leq S_{10}+\frac{1}{10}$}
\item []
\item []<5->Which gives an estimate for $S=\sum_{n=1}^{\infty}\frac{1}{n^{2}}$,
$1.64068\leq S\leq1.64977$. 
\end{itemize}
\end{frame}

\begin{frame}{Homework}

\textbf{HW: }12, 14, 20, 24, 29
\end{frame}

\end{document}
